
% Page 55

La paix est donc officiellement établie par Henri IV, mais la guerre de Rohan qui suivra quelques années après aggravera la situation au point que, 140 ans après, les États albigeois de 1740 mentionneront la détresse des habitants d'Alban et porteront secours à la population. Ces gestes secourables se renouvelleront en 1763 et 1783 : deux siècles d'expiation !

Faut-il mettre au compte de ces années mouvementées le calvaire de pierre qui se trouve actuellement dans l'église et qui fut classé comme étant du XVI\textsuperscript{ème} siècle : ce qui contredit -- sans offenser personne -- l'indication donnée par le panneau qui le signale aux visiteurs occasionnels. La scène centrale du piédroit, la plus basse, y représente en effet Véronique tenant le linge sur lequel s'est imprimée la face du Sauveur, après que la pieuse israélite eut essuyé le visage du Seigneur sur le chemin du calvaire. Or, cette représentation n'apparaît pas dans l'iconographie chrétienne avant le XV\textsuperscript{ème} siècle. De plus, la rose des Tudor qui figure en ornementation juste au-dessous du tableau supérieur représentant la crucifixion, se retrouve exactement semblable dans une pierre gravée et en remploi sur la façade actuelle de l'église de Miolles. Or, la date qui accompagne ces diverses sculptures de pierre est 1608. Ceci donne donc à penser que, tout en ayant une facture très archaïque, rappelant davantage le roman, notre calvaire pourrait bien être seulement l'œuvre d'un artisan attardé dans ses procédés de composition sculpturale. Cela n'enlève rien à la valeur du monument.
